\documentclass[11pt]{article}
\usepackage{fontspec}
\setmainfont{Latin Modern Roman}
\usepackage{amsmath,amssymb,mathtools,unicode-math}
\setmathfont{Latin Modern Math}
\usepackage{xurl,hyperref}
\hypersetup{hidelinks,pdftitle={Problem statement}}
\newcommand{\sref}[2]{\hyperlink{source:#1:#2}{[S#2]}}
\newcommand{\cataloguescope}[1]{\paragraph{Mathematical question.}#1}
\begin{document}
\section{Circular flow number of regular class 1 graphs}\label{statement}

\subsection{Definitions and mathematical statement}

\paragraph{Shared notation.}$\mathbb N=\{1,2,\ldots\}$, and $\mathbb Z,\mathbb Q,\mathbb R,\mathbb C$ denote the integers, rationals, reals and complex numbers. Any other conventions are specified locally.


Graphs are finite and loopless and may have parallel edges. A nowhere-zero real $r$-flow consists of an orientation and real numbers $\phi(e)$ such that $1\le|\phi(e)|\le r-1$ and, at every vertex, the sum on outgoing edges equals the sum on incoming edges. The circular flow number $F_c(G)$ is the infimum of such real $r\ge2$, with value $+\infty$ if none exist. For $X\subseteq V(G)$, $\delta_G(X)$ is the set of edges with exactly one endpoint in $X$. A graph is class $1$ if its chromatic index equals its maximum degree.
\paragraph{Statement recorded in the supplied source.}
For every integer $t\ge1$ and every $(2t+1)$-regular class-$1$ graph $G$, is $F_c(G)\le2+2/t$? \sref{statement}{2}
\subsection{Short English statement}
Does an optimal edge coloring of an odd regular graph force the proposed circular-flow bound?
\subsection{Sources}
\begin{description}
\item[\hypertarget{source:statement:1}{S1}] ULAM.AI and the UnsolvedMath Contributors, UnsolvedMath: A Curated Collection of Open Mathematics Problems, catalogue entry. CC BY 4.0. \url{https://huggingface.co/datasets/ulamai/UnsolvedMath}.
\item[\hypertarget{source:statement:2}{S2}] Open Problem Garden, Circular flow number of regular class 1 graphs, original node 59994, statement and mathematical discussion. \url{https://www.openproblemgarden.org/op/circular_flow_number_of_regular_class_1_graphs}.
\end{description}
\label{end:statement}
\end{document}
